\documentclass{acm_proc_article-sp}

\usepackage[utf8]{inputenc}
\usepackage[a4paper,margin=1.2cm,bottom=0.5cm,top=1cm]{geometry}
\usepackage{graphicx}
\usepackage{wrapfig}

\newcommand{\fig}[4][htb]{
  \begin{figure}[#1]
    {\centering{\includegraphics[width=#4]{fig/#2}}\par}
  % \caption{#3}
    \label{fig:#2}
  \end{figure}
}

\newcommand{\ffig}[4][r]{
  \begin{wrapfigure}{#1}{#4}
    {\centering{\includegraphics[width=#4]{fig/#2}}\par}
  % \caption{#3}
    \label{fig:#2}
  \end{wrapfigure}
}

\newcommand{\tab}[4][htb]{
  \begin{table}[h]
    {\centering\footnotesize\textsf{\input{fig/#2.tab}}\par}
%    \caption{#3}
    \label{tab:#2}
  \end{table}
}

\begin{document}

\title{
  \begin{tabular*}{\textwidth}{@{\extracolsep{\fill}}p{2cm}cp{2cm}}
    \raisebox{-9mm}[0pt][0pt]{\includegraphics[height=1.5cm]{fig/ufsc}}
    &\textbf{\large Federal University of Santa Catarina}
    &\raisebox{-8mm}[0pt][0pt]{\includegraphics[height=1cm]{fig/logolisha}}\\
    &\textbf{\large Software/Hardware Integration Lab}&\\
    \noalign{\smallskip}
    &\texttt{http://www.lisha.ufsc.br/}&\\
    \noalign{\bigskip}
    \noalign{\smallskip}
    \multicolumn{3}{c}{\centering\textbf{\LARGE Smart Cities and the Internet of Things}}\\
    \noalign{\bigskip}
  \end{tabular*}
}

\maketitle

\section*{Smart Cities}

The way we live has been changing quickly under the influence of the information era technologies. The cities in which we live are the next fated to follow the revolution. \emph{Connectivity} is becoming ubiquitous to modern cities and the foretold massive deployment of \emph{sensors} will soon enable us to interact with our urban environment using the same tools we now use to access information. Policy and decision makers will be supported by data of unprecedented magnitude.
\fig{smartcity}{nocaption}{.9\columnwidth}

\section*{Smart Cities on the SmartGrid}

The way we produce and distribute electricity has remained mostly
unchanged since Telsa's time. However, the recent recognition that
natural resources are being depleted faster than the planet can
replenish them drove us toward more \emph{sustainable} strategies. The
\emph{SmartGrid} combines ICT and EE technology to improve the power
grid efficiency and to enable distributed generation, the key to the
integration of \emph{clean energy} sources.  The underlying technology
shares many aspects with that for \emph{Smart Cities}.  Addressing both
challenges at once will certainly optimize the endeavor.

\fig{smartgrid}{nocaption}{.9\columnwidth}

\ffig{broadband}{nocaption}{2cm}
\section*{Broadband Internet}

\emph{Broadband internet access} is another keystone for \emph{Smart
  Cities}, providing the pathways for data acquisition and service
supply. In particular, digital media services are being quickly
incorporated into both the Internet of People and the \emph{Internet of
  Things}: IPTV, VoD, monitoring, and surveillance are just the
triggering apps of an unforeseeable list of services.


\section*{Internet of Things}

\ffig{iot}{nocaption}{3.5cm} 

The \emph{Internet of Things} will allow objects to interact with each
other just the way people now interact over the Internet.  Devices,
appliances, machines, houses, and buildings are quickly incorporating
the technological offspring of Wireless Sensor Networks to become
\emph{Smart Things} that are able to adjust their behavior in order to
spare resources and to improve usability.

\section*{The Role of LISHA} 

UFSC's Software/Hardware Integration Lab was founded in 1984 to promote
research in the frontiers between hardware and software. Since then, the
lab has been carrying on research activities that now makes it a key
partner on the pursue of \emph{Smart Cities}.

\ffig{adesd_tools}{nocaption}{4.6cm}
\subsection*{ADESD and EPOS}

LISHA's major research back-plane, the \emph{Application-driven Embedded
  System Design}~(ADESD) methodology guides the development of embedded
systems as aggregate of adaptable HW/SW components. EPOS, its best show
case, has been ported to several architectures and now supports
innumerable applications.

\ffig[l]{emote}{nocaption}{2cm}
\subsection*{EPOSMote} 

EPOSMote is an \emph{IoT mote} specifically designed to run EPOS. The
\emph{free hardware} platform is now being used by several universities.

\ffig[l]{cia2}{nocaption}{2cm}
\subsection*{CIA$^2$} 

LISHA is currently co-leading the CIA$^2$ research network, which
encompasses 18 Brazilian universities on the quest for Smart Cities
infrastructure and applications.
 
\ffig[l]{altatv}{nocaption}{2cm}
\subsection*{ALTATV} 

LISHA led the research network that developed the ALTATV Open Free
Scalable Access Terminal Architecture: the link between the \emph{Smart
  Home} and the Internet.

\ffig[l]{eboat}{nocaption}{2cm}
\subsection*{Photovoltaics Lab} 

A joint venture with UFSC's Photovoltaics Lab established two key
deployment cases for LISHA's technology: \emph{Smart Buildings} and
\emph{Electric Vehicles}.

\begin{table}[h]
  \begin{tabular*}{\columnwidth}{@{\extracolsep{\fill}}l}
    \hline
\\[-8pt]
    Prof.~Dr.~Antônio Augusto Fröhlich\\
    UFSC/CTC/LISHA\\
    88040-900 Florianópolis - SC - BRAZIL\\
    \texttt{http://www.lisha.ufsc.br/~~~+55 48 3721-9516}
  \end{tabular*}
\end{table}

\end{document}

